\begin{proof}[Proof of \cref{obj:prop:fixed-d-effective-sample-reduction}]
Write
\[
N(t)=n(t)q(t)>0,
\qquad
\ell(t)=\log(e+N(t)).
\]
Define the restricted model class and its minimax risk by
\[
\mathcal M^{(0)}_{n,d,q}
:=
\left\{
P\in\mathcal M_{n,d,q}:
P(S(0)=0,S(1)=0)=1
\right\},
\qquad
\mathfrak R^{(0)}_{n,d,q}
:=
\inf_{\mathsf T\in\mathcal T_n}
\sup_{P\in\mathcal M^{(0)}_{n,d,q}}
E_P[(\mathsf T-\tau(P))^2].
\]
Here the estimator class and joint sample-and-kernel expectation are those of
\cref{obj:def:minimax-risk}.

\begin{enumerate}
\item Since \(N(t)\to\infty\), eventually
\[
N(t)\ge \max\{2,d^2\}.
\]
Also \(\ell(t)\ge1\), because \(e+N(t)\ge e\). Consequently, on this eventual range,
\[
\left[\frac{d}{N(t)\ell(t)}\right]^2
=
\frac{d^2}{N(t)^2\ell(t)^2}
\le
\frac{1}{N(t)\ell(t)^2}
\le
N(t)^{-1}.
\]
The expression in \cref{obj:def:frontier-rate} is nonnegative for every \(t\).
Eventually \(N(t)^{-1}\le1\), and both entries of its minimum are at least
\(N(t)^{-1}\). Thus
\[
N(t)^{-1}
\le
r_{n(t),d,q(t)}
=
\min\left\{
1,\,
N(t)^{-1}+
\left[\frac{d}{N(t)\ell(t)}\right]^2
\right\}
\le
2N(t)^{-1}.
\]
These inequalities prove the first order comparison.
% lean: CausalSmith.Stat.MarRareqLogfrontier.fixed_d_frontier_rate_theta

\item Let
\[
0<\underline c_{\mathrm{match}}
\le \overline C_{\mathrm{match}}<\infty
\]
be the constants supplied by \cref{obj:thm:matched-minimax-frontier}.
Its hypotheses hold at every \(t\): the sample and alphabet sizes are positive,
\(0<q(t)\le1\), and the rare-arrival restriction is assumed. Hence
\[
\underline c_{\mathrm{match}}\,r_{n(t),d,q(t)}
\le
\mathfrak R_{n(t),d,q(t)}
\le
\overline C_{\mathrm{match}}\,r_{n(t),d,q(t)}
\qquad(t\in\mathbb N).
\]
In particular, both quantities are nonnegative, and the lower inequality also gives
\[
r_{n(t),d,q(t)}
\le
\underline c_{\mathrm{match}}^{-1}
\mathfrak R_{n(t),d,q(t)}.
\]
Combining these bounds with the eventual frontier bounds from the preceding step yields
\[
\underline c_{\mathrm{match}}N(t)^{-1}
\le
\mathfrak R_{n(t),d,q(t)}
\le
2\overline C_{\mathrm{match}}N(t)^{-1}.
\]
% lean: CausalSmith.Stat.MarRareqLogfrontier.matched_minimax_frontier

\item We next obtain a pointwise upper bound for the restricted risk.
For each \(t\), use the estimator of \cref{obj:def:mixed-count-estimator}
with parameters \(n(t),d,q(t)\). Its auxiliary output is zero or a value clipped
to \([-1,1]\), so its conditional average satisfies
\[
-1
=
E_\omega[-1\mid\mathbf O_{n(t)}]
\le
\widehat\tau^{\mathrm{mix}}_{n(t),d,q(t)}(\mathbf O_{n(t)})
\le
E_\omega[1\mid\mathbf O_{n(t)}]
=
1.
\]
The average is measurable because it is a finite sum over the admissible prefix
lengths and assignments. Therefore the deterministic kernel
\[
\mathsf T^{\mathrm{mix}}_t(o,dh)
:=
\operatorname{Dirac}_{\widehat\tau^{\mathrm{mix}}_{n(t),d,q(t)}(o)}(dh)
\]
belongs to \(\mathcal T_{n(t)}\). Its joint risk equals its deterministic risk:
\[
E_P[(\mathsf T^{\mathrm{mix}}_t-\tau(P))^2]
=
E_P\!\left[
\left(
\widehat\tau^{\mathrm{mix}}_{n(t),d,q(t)}(\mathbf O_{n(t)})
-\tau(P)
\right)^2
\right].
\]
Let
\[
0<\overline C_{\mathrm{mix}}<\infty
\]
be supplied by \cref{obj:thm:mixed-count-upper}. For every
\(P\in\mathcal M^{(0)}_{n(t),d,q(t)}\), membership in the full model class
allows us to apply its risk-envelope bound:
\[
E_P[(\mathsf T^{\mathrm{mix}}_t-\tau(P))^2]
\le
\mathfrak U_{n(t),d,q(t)}
\le
\overline C_{\mathrm{mix}}\,r_{n(t),d,q(t)}.
\]
All squared risks are nonnegative, so the minimax value is bounded above by
the worst-case risk of this admissible kernel. If the restricted class is
nonempty, taking its supremum gives
\begin{equation}\label{eq:fixed-d-restricted-risk-upper}
\mathfrak R^{(0)}_{n(t),d,q(t)}
\le
\sup_{P\in\mathcal M^{(0)}_{n(t),d,q(t)}}
E_P[(\mathsf T^{\mathrm{mix}}_t-\tau(P))^2]
\le
\overline C_{\mathrm{mix}}\,r_{n(t),d,q(t)}.
\end{equation}
If the restricted class is empty, its real-valued worst-case risk is zero
under the empty-supremum convention. The same inequality follows from
\[
0\le \overline C_{\mathrm{mix}}\,r_{n(t),d,q(t)}.
\]
Thus the upper bound holds for every \(t\).
% lean: CausalSmith.Stat.MarRareqLogfrontier.mixed_count_upper
% lean: CausalSmith.Stat.MarRareqLogfrontier.fixed_d_effective_sample_reduction

\item To derive the restricted lower bound, fix \(t\). For
\(\mu_{\mathrm{test}}\in[1/4,3/4]\), define \(P_{\mu_{\mathrm{test}}}^{(t)}\) by
\[
X=0,\qquad S(0)=S(1)=Y(0)=0,
\qquad
A\sim\operatorname{Bernoulli}(1/2),\quad
Y(1)\sim\operatorname{Bernoulli}(\mu_{\mathrm{test}}),\quad
R\sim\operatorname{Bernoulli}(q(t)),
\]
where \(A,Y(1),R\) are mutually independent, and set
\(S=S(A)\), \(Y=Y(A)\). The category \(0\) exists because \(d\ge1\).
Take the observed sample independently from the induced observed-record law.

This construction satisfies randomization and consistency, and independence
of \(R\) from \((X,A,S,Y)\) implies the required conditional arrival independence.
Its occupied cells satisfy
\[
p_{a,0,0}(P_{\mu_{\mathrm{test}}}^{(t)})=\frac12,
\qquad
\rho_{a,0,0}(P_{\mu_{\mathrm{test}}}^{(t)})=q(t)
\qquad(a\in\{0,1\}),
\]
while every other cell has zero membership probability. Together with the
assumed parameter restrictions and rare-arrival slice, these facts establish
\[
P_{\mu_{\mathrm{test}}}^{(t)}\in\mathcal M^{(0)}_{n(t),d,q(t)},
\qquad
\tau(P_{\mu_{\mathrm{test}}}^{(t)})
=
E_{P_{\mu_{\mathrm{test}}}^{(t)}}[Y(1)-Y(0)]
=
\mu_{\mathrm{test}}.
\]

Define the separation and the two outcome means by
\[
\delta(t)
:=
\min\left\{\frac14,\frac{1}{4\sqrt{N(t)}}\right\},
\qquad
u_0(t):=\frac12,
\qquad
u_1(t):=\frac12+\delta(t).
\]
Since \(N(t)>0\),
\[
0<\delta(t)\le\frac14,
\qquad
u_0(t),u_1(t)\in[1/4,3/4].
\]
For \(i\in\{0,1\}\), define the single-record and sample laws by
\[
Q_i^{\mathrm{obs}}(t)
:=
\mathcal L_{P_{u_i(t)}^{(t)}}(O_1),
\qquad
Q_i^{\mathrm{sam}}(t)
:=
\mathcal L_{P_{u_i(t)}^{(t)}}(\mathbf O_{n(t)})
=
\bigl(Q_i^{\mathrm{obs}}(t)\bigr)^{\otimes n(t)}.
\]

For finite probability laws \(Q_1\ll Q_0\), use the divergence and
total-variation conventions
\[
\chi^2(Q_1\Vert Q_0)
:=
\sum_{o:\,Q_0(\{o\})>0}
\frac{\bigl(Q_1(\{o\})-Q_0(\{o\})\bigr)^2}{Q_0(\{o\})},
\qquad
\operatorname{TV}(Q_1,Q_0)
:=
\sup_{\mathcal A}|Q_1(\mathcal A)-Q_0(\mathcal A)|
=
\frac12\sum_o|Q_1(\{o\})-Q_0(\{o\})|.
\]
Under \(P_{\mu_{\mathrm{test}}}^{(t)}\), every observed record has \(X=S=0\). The remaining
coordinates, with arrived outcome \(RY=R\,Y\), have probabilities
\[
\begin{array}{c|ccccc}
(A,R,RY)
&(0,0,0)&(1,0,0)&(0,1,0)&(1,1,0)&(1,1,1)\\ \hline
\text{Probability}
&\dfrac{1-q(t)}2
&\dfrac{1-q(t)}2
&\dfrac{q(t)}2
&\dfrac{q(t)(1-\mu_{\mathrm{test}})}2
&\dfrac{q(t)\mu_{\mathrm{test}}}2
\end{array}
\]
and all other records have probability zero. At \(q(t)=1\), the first two
atoms have zero probability under both laws and are omitted from the
divergence sum. Thus \(Q_1^{\mathrm{obs}}(t)\ll Q_0^{\mathrm{obs}}(t)\), and summing
the atom probabilities gives
\[
\begin{aligned}
1+\chi^2\!\left(Q_1^{\mathrm{obs}}(t)\Vert Q_0^{\mathrm{obs}}(t)\right)
&=
1-q(t)+\frac{q(t)}2
+q(t)\bigl((1-u_1(t))^2+u_1(t)^2\bigr)\\
&=
1+2q(t)\delta(t)^2.
\end{aligned}
\]
The definition of \(\delta(t)\) gives
\[
16N(t)\delta(t)^2\le1,
\qquad
\chi^2\!\left(Q_1^{\mathrm{obs}}(t)\Vert Q_0^{\mathrm{obs}}(t)\right)
=
2q(t)\delta(t)^2
\le
\frac{1}{8n(t)}.
\]
Factoring the finite product sums and using \(1+x\le e^x\) now yields
\[
\begin{aligned}
1+\chi^2\!\left(Q_1^{\mathrm{sam}}(t)\Vert Q_0^{\mathrm{sam}}(t)\right)
&=
\left(
1+\chi^2\!\left(Q_1^{\mathrm{obs}}(t)\Vert Q_0^{\mathrm{obs}}(t)\right)
\right)^{n(t)}\\
&\le
\exp\!\left(
n(t)\chi^2\!\left(Q_1^{\mathrm{obs}}(t)\Vert Q_0^{\mathrm{obs}}(t)\right)
\right)
\le e^{1/8}.
\end{aligned}
\]
Since
\[
e^{1/2}e^{1/2}=e<3,
\qquad
e^{1/8}\le e^{1/2}<\sqrt3<2,
\]
the sample divergence is less than \(1\). Cauchy--Schwarz applied to the
finite total-variation sum therefore gives
\[
\operatorname{TV}\!\left(Q_1^{\mathrm{sam}}(t),Q_0^{\mathrm{sam}}(t)\right)
\le
\frac12
\sqrt{\chi^2\!\left(Q_1^{\mathrm{sam}}(t)\Vert Q_0^{\mathrm{sam}}(t)\right)}
\le\frac12.
\]
Consequently, every measurable sample event \(\mathcal A\) satisfies
\[
Q_1^{\mathrm{sam}}(t)(\mathcal A^c)
+
Q_0^{\mathrm{sam}}(t)(\mathcal A)
\ge
1-\operatorname{TV}\!\left(Q_1^{\mathrm{sam}}(t),Q_0^{\mathrm{sam}}(t)\right)
\ge\frac12.
\]
% lean: CausalSmith.Stat.MarRareqLogfrontier.oneCell_product_testing_half

\item Let \(\mathsf T\in\mathcal T_{n(t)}\) be arbitrary, and define its
barycenter and the relevant sample events by
\[
b_{\mathsf T}(o):=\int_{[-1,1]}h\,\mathsf T(o,dh),
\]
\[
\mathcal E_{\mathsf T,i}
:=
\left\{o:\frac{\delta(t)}2\le
|b_{\mathsf T}(o)-u_i(t)|\right\}
\quad(i\in\{0,1\}),
\qquad
\mathcal A_{\mathsf T}
:=
\left\{o:
|b_{\mathsf T}(o)-u_1(t)|
\le
|b_{\mathsf T}(o)-u_0(t)|
\right\}.
\]
These events are measurable. The triangle inequality gives, at every sample,
\[
\delta(t)
=
|u_1(t)-u_0(t)|
\le
|b_{\mathsf T}(o)-u_0(t)|
+
|b_{\mathsf T}(o)-u_1(t)|.
\]
It follows that
\[
\mathcal A_{\mathsf T}\subseteq\mathcal E_{\mathsf T,0},
\qquad
\mathcal A_{\mathsf T}^{c}\subseteq\mathcal E_{\mathsf T,1}.
\]
Applying the testing bound from the preceding step and then these inclusions,
\[
\frac12
\le
Q_1^{\mathrm{sam}}(t)(\mathcal A_{\mathsf T}^{c})
+
Q_0^{\mathrm{sam}}(t)(\mathcal A_{\mathsf T})
\le
Q_1^{\mathrm{sam}}(t)(\mathcal E_{\mathsf T,1})
+
Q_0^{\mathrm{sam}}(t)(\mathcal E_{\mathsf T,0}).
\]
Thus
\[
\frac14
\le
\max_{i\in\{0,1\}}
Q_i^{\mathrm{sam}}(t)(\mathcal E_{\mathsf T,i}).
\]
On \(\mathcal E_{\mathsf T,i}\), the barycenter's squared error is at least
\((\delta(t)/2)^2\). The barycenter inequality of
\cref{obj:lem:randomized-kernel-barycenter} therefore supplies
\[
\left(\frac{\delta(t)}2\right)^2
Q_i^{\mathrm{sam}}(t)(\mathcal E_{\mathsf T,i})
\le
\int
(b_{\mathsf T}(o)-u_i(t))^2\,Q_i^{\mathrm{sam}}(t)(do)
\le
E_{P_{u_i(t)}^{(t)}}[(\mathsf T-u_i(t))^2].
\]
Whether the probability for \(i=0\) is at most the probability for \(i=1\),
or conversely, multiplication by the nonnegative factor
\((\delta(t)/2)^2\) preserves their order. Hence
\[
\begin{aligned}
\frac{\delta(t)^2}{16}
&\le
\left(\frac{\delta(t)}2\right)^2
\max_{i\in\{0,1\}}
Q_i^{\mathrm{sam}}(t)(\mathcal E_{\mathsf T,i})\\
&=
\max_{i\in\{0,1\}}
\left\{
\left(\frac{\delta(t)}2\right)^2
Q_i^{\mathrm{sam}}(t)(\mathcal E_{\mathsf T,i})
\right\}\\
&\le
\max_{i\in\{0,1\}}
E_{P_{u_i(t)}^{(t)}}[(\mathsf T-u_i(t))^2].
\end{aligned}
\]
% lean: CausalSmith.Stat.MarRareqLogfrontier.complete_arrival_two_point_decision_reduction

\item The separation also has the required lower calibration.
If \(1/4\le1/(4\sqrt{N(t)})\), then
\[
\delta(t)^2=\frac1{16}
\ge
\frac1{16}\min\{1,N(t)^{-1}\}.
\]
In the other case,
\[
\delta(t)^2
=
\frac{1}{16N(t)}
\ge
\frac1{16}\min\{1,N(t)^{-1}\}.
\]
Combining either inequality with the preceding risk bound gives, for every
\(\mathsf T\in\mathcal T_{n(t)}\),
\[
\frac1{256}\min\{1,N(t)^{-1}\}
\le
\frac{\delta(t)^2}{16}
\le
\max_{i\in\{0,1\}}
E_{P_{u_i(t)}^{(t)}}[(\mathsf T-u_i(t))^2].
\]

The estimator class is nonempty: the constant kernel
\[
\mathsf T^{0}(o,dh):=\operatorname{Dirac}_{0}(dh)
\]
is admissible. Moreover, binary potential outcomes imply
\(\tau(P)\in[-1,1]\), and the kernel outputs lie in \([-1,1]\), so
\[
0\le E_P[(\mathsf T-\tau(P))^2]\le4
\]
for every full-data law \(P\) and admissible \(\mathsf T\).
Both testing laws belong to the full model and to its restricted subclass.
Their maximum risk therefore bounds each corresponding worst-case risk
from below; the uniform bound by \(4\) ensures these suprema are finite.
Taking the infimum over estimators gives
\[
\frac1{256}\min\{1,N(t)^{-1}\}
\le
\mathfrak R_{n(t),d,q(t)},
\qquad
\frac1{256}\min\{1,N(t)^{-1}\}
\le
\mathfrak R^{(0)}_{n(t),d,q(t)}.
\]
% lean: CausalSmith.Stat.MarRareqLogfrontier.oneCell_minimax_risk_floors

\item The last lower bound establishes nonnegativity of the restricted risk
for every \(t\). Its upper bound in \cref{eq:fixed-d-restricted-risk-upper}, combined with the eventual
frontier bound, gives
\[
0\le
\mathfrak R^{(0)}_{n(t),d,q(t)}
\le
2\overline C_{\mathrm{mix}}N(t)^{-1}
\qquad\text{eventually}.
\]
Also \(N(t)\to\infty\) implies \(N(t)\ge1\) eventually, so
\[
\min\{1,N(t)^{-1}\}=N(t)^{-1}.
\]
The lower bound consequently becomes
\[
\frac1{256}N(t)^{-1}
\le
\mathfrak R^{(0)}_{n(t),d,q(t)},
\qquad\text{equivalently}\qquad
N(t)^{-1}\le256\,\mathfrak R^{(0)}_{n(t),d,q(t)}.
\]
Together these inequalities prove the restricted-risk order comparison,
and hence all three assertions.
% lean: CausalSmith.Stat.MarRareqLogfrontier.fixed_d_effective_sample_reduction
\end{enumerate}
\end{proof}
